\documentclass[10pt]{article}

\usepackage[utf8]{inputenc}

\usepackage[T1]{fontenc}

\usepackage{amsmath}

\usepackage{amsfonts}

\usepackage{amssymb}

\usepackage[version=4]{mhchem}

\usepackage{stmaryrd}

\usepackage{graphicx}

\usepackage[export]{adjustbox}

\graphicspath{ {./images/} }



\begin{document}

где $z^{\prime}-$ ось, проходящая через центр масс тела, а $d-$ расстояние между осями.



М. и. относительно любой проходящей через начало координат $O$ оси $O l$ с направляющими косинусами $\alpha, \beta, \gamma$ находится по формуле:



\[

I_{O l}=I_{x} \alpha^{2}-I_{y} \beta^{2}-I_{z} \gamma^{2}-2 I_{x y} \alpha \beta-2 I_{y z} \beta \gamma-2 I_{z x} \gamma \alpha

\]



Зная шесть величин $I_{x}, I_{y}, I_{z}, I_{x y}, I_{y z}, I_{z x}$, можно последовательно, используя формулы (4) и (3), вычислить всю совокупность М. и. тела относительно любых осей. Эти шесть величин определяют так наз. те нзор инерции тела. Через каждую точку тела можно провести три такие взаимно перпендикулярные оси, называемые гл а в ны м и ос ям и и н е р ци и, для к-рых



\[

I_{x y}=I_{y z}=I_{z x}=0 .

\]



Тогда, зная главные оси инерции и М. и. относительно этих осей, можно определить М. и. тела относительно любой оси.



Лит.: [1] Ф а в о р и н М. В., Моменты инерции тел. Справочник, М., 1970; [2] Гернет М. М., Ратобыльский В. Ф., Определение моментов инерции, М., $1969 . \quad$ С. М. Тарг. МОМЕНТ КОЛИЧЕСТВА ДВИХЕНИЯ, кине тиче ский момент, угловой момент,- одна из мер механического движения материальной точки или системы. Особенно важную роль М. К. д. играет при изучении вращательного движения. Как и для момента силы, различают М. к. Д. относительно центра (точки) и относительно оси.



Для вычисления М.К. д. К материальной точки относительно центра $O$ или оси $z$ справедливы все формулы, приведенные для вычисления момента силы, если в них заменить вектор $\boldsymbol{F}$ вектором количества движения $\boldsymbol{m} \boldsymbol{v}$. Таким образом,



\[

\boldsymbol{K}_{O}=|\boldsymbol{r} \cdot \boldsymbol{m v}|

\]



где $\boldsymbol{r}-$ радиус-вектор движущейся точки, проведенный из центра $O$, а $K_{z}$ равняется проекции вектора $K_{O}$ на ось $z$, проходящую через точку $О$. Изменение М. к. д. точки происходит под действием момента $\boldsymbol{m}_{O}(\boldsymbol{F})$ приложенной силы и определяется теоремой об изменении М.к.д., выражаемой уравнением



\[

d \boldsymbol{K}_{O} / d t=\boldsymbol{m}_{O}(\boldsymbol{F})

\]



Когда $\boldsymbol{m}_{O}(\boldsymbol{F})=0$, что, напр., имеет место для центральных сил, движение точки подчиняется площадей закону. Этот результат важен для небесной механики, теории движения искусственных спутников Земли, космич. летательных аппаратов и др.



Главный М.к.д. (или кинетический момент) механич. системы относительно центра $O$ или оси $z$ равен соответственно геометрич. или алгебраич. сумме М.к.д. всех точек системы относительно того же центра или оси, то есть



\[

K_{O}=\sum K_{O i}, \quad K_{z}=\sum K_{z i}

\]



Вектор $K_{O}$ может быть определен его проекциями $K_{x}, K_{y}, K_{z}$ на координатные оси. Для тела, к-рое вращается вокруг неподвижной оси $z$ с угловой скоростью $\omega$,



\[

K_{x}=-I_{x z} \omega, K_{y}=-I_{y z} \omega, K_{z}=I_{z} \omega,

\]



где $I_{z}$ - осевой, а $I_{x z}, I_{y z}$ - центробежные моменты инерции. Если ось $z$ является главной осью инерции для начала координат $O$, то $\boldsymbol{K}_{O}=I_{z} \omega$.



Изменение главного М.к. д. системы происходит под действием только внешних сил и зависит от их главного момента $M_{O}^{e}$. Эта зависимость определяется теоремой об изменении главного М. к. д. системы, выражаемой уравнением



\[

d K_{O} / d t=M_{O}^{e}

\]



Аналогичным уравнением связаны моменты $K_{z}$ и $M_{z}^{e}$ относительно оси $z$. Если $M_{O}^{e}=0$ или $M_{z}^{e}=0$, то соответственно $K_{O}$ или $K_{z}$ будут величинами постоянными, то есть имеет место закон сохранения М. к. Д. (см. Сохранения законы). Таким образом, внутренние силы не могут изменить М. К. Д. системы, но М.к. д. отдельных частей системы или угловые скорости под действием этих сил могут изменяться. Напр., у вращающегося вокруг вертикальной оси $z$ фигуриста (или балерины) величина $K_{z}=I_{z} \omega$ будет постоянной, так как практически $M_{z}^{e}=0$. Но, изменяя движением рук или ног значение момента инерции $I_{z}$, фигурист может изменять угловую скорость $\omega$. Понятие о М. к. д. широко используется в динамике твердогочтела, особенно в теории гироскопа.



Размерность М.к. д.- $L^{2} M T^{-1}$. М. к. д. обладают также электромагнитные, гравитационные и другие физич. поля. Большинству элементарных частиц присущ собственный, внутренний М. К. д.- спин. Большое значение М. К. д. имеет в квантовой механике.



См. также Орбитальный момент. МОМЕНт Силы - величина, характеризующая вращательный эффект силы при действии ее на твердое тело; одно из основных понятий механики. Различают М.с. относительно центра (точки) и относительно оси.



М. с. относительно центра $O-$ величина векторная. Его модуль $M_{O}=F h$, где $F$ - модуль силы, а $h$ - плечо, то есть длина перпендикуляра, опущенного из $O$ на линию действия силы (см. рис.); направлен вектор $\boldsymbol{M}_{O}$ перпендикулярно плоскости, проходящей через центр $O$ и силу $F_{\text {в сторо- }}$



Вектор $\boldsymbol{M}_{O}-$ момент силы $F$ относительно центра $O$.



\begin{center}

\includegraphics[max width=\textwidth]{2023_07_08_0bc501b6b92a1608faccg-1}

\end{center}



ну, откуда поворот, вызываемый силой, виден против хода часовой стрелки (в правой системе координат). С помощью векторного произведения М. с. выражается равенством



\[

\boldsymbol{M}_{O}=[\boldsymbol{r} \cdot \boldsymbol{F}],

\]



где $\boldsymbol{r}-$ радиус-вектор, проведенный из $O$ в точку приложения силы. Размерность М.с. $L^{2} M T^{2}$.



М. с. относительно оси - величина алгебраическая, равная проекции на эту ось М. с. относительно любой точки $O$ оси или же численной величине момента проекции $F_{x y}$ силы $\boldsymbol{F}$ на плоскость $x y$, перпендикулярную оси $z$, взятого относительно точки пересечения оси с плоскостью. То есть



\[

M_{z}=M_{O} \cos \gamma=F_{x y} h_{1} \text { или } M_{z}=-F_{x y} h_{1} .

\]



Знак минус (последнее выражение) берется, когда поворот силы $\boldsymbol{F}$ с положительного направления оси $z$ виден по ходу часовой стрелки (тоже в правой системе). М. с. относительно осей $x, y, z$ могут также вычисляться по формулам:



\[

\begin{aligned}

& M_{x}=y F_{z}-z F_{y}, \\

& M_{y}=z F_{x}-x F_{z}, \\

& M_{z}=x F_{y}-y F_{x},

\end{aligned}

\]



где $F_{x}, F_{y}, F_{z}$ - проекции силы $F_{\text {на оси; }} x, y, z-$ координаты точки $A$ приложения силы.



Если система сил имеет равнодействующую, то ее момент вычисляется по Вариньона теореме.



C. M. Тарг.



MOMEHTA ОТОБРАЖEHИЕ - отображение симплектического многообразия $(M, \omega)$ с пуассоновым действием





\end{document}